\section{Firmware}
\label{chap:imp_firmware}

With all hardware components modeled and integrated into the QEMU emulator, only 
the firmware for loading onto the Netduino Plus 2 board and controlling the LIS3DH 
accelerometer remains. Although Seeed-Studio provides an open-source library for 
managing the LIS3DH accelerometer\footnote{\url{https://github.com/Seeed-Studio/Seeed_Arduino_LIS3DHTR}}, 
this library has the disadvantage of being designed for Arduino or Raspberry Pi 
boards, not for STM32-based platforms. Therefore, based on the Seeed-Studio library, 
a custom C library was developed to manage the LIS3DH accelerometer.

This library enables device configuration and reading of acceleration and 
temperature data through the I$^2$C bus. It is designed for use with STM32F4 series 
microcontrollers, specifically the STM32F405, utilizing STMicroelectronics' 
Hardware Abstraction Layer (HAL) as its foundation. The library comprises three main files:

\begin{enumerate}
    \item \textbf{\texttt{lis3dh\_types.h}}: Contains data type definitions, 
    enumerations, and structures necessary for the sensor interface.
    
    \item \textbf{\texttt{lis3dh\_driver.h}}: Defines constants, macros, and public 
    function prototypes for the driver.
    
    \item \textbf{\texttt{lis3dh\_driver.c}}: Implements the configuration and 
    communication logic with the LIS3DH sensor.
\end{enumerate}

\subsection{Data Structures}

To encapsulate accelerometer information, the following structures were defined in 
\texttt{lis3dh\_types.h}:

\begin{itemize}
    \item \textbf{\texttt{lis3dh\_config\_t}}: Stores device configuration, including 
    measurement scale (2g, 4g, 8g, 16g), operating mode (high resolution, normal, 
    low power), output data rate (ODR), and enablement of temperature sensor and 
    high-pass filter.
    
    \item \textbf{\texttt{lis3dh\_data\_t}}: Contains sensor readings, i.e., 
    accelerations on three axes (in mg) and temperature (in degrees Celsius).
    
    \item \textbf{\texttt{lis3dh\_t}}: Main structure grouping the I$^2$C handler, 
    device address, configuration, sensitivity, and data.
\end{itemize}

\subsection{Implemented Functionality}

The library provides the following functions:

\subsubsection{Initialization and Configuration}

\begin{itemize}
    \item \texttt{lis3dh\_init}: Initializes the accelerometer structure, verifies 
    I$^2$C communication, and configures default or user-provided parameters.
    
    \item \texttt{lis3dh\_configuration}: Applies specific device configuration, 
    setting operating mode, scale, sampling rate, and enabling axes.
\end{itemize}

\subsubsection{Interrupt Configuration}

\begin{itemize}
    \item \texttt{lis3dh\_enable\_drdy\_interrupt}: Enables Data-Ready (DRDY) 
    interrupt on the accelerometer's INT1 pin, notifying the microcontroller when 
    new data is available.
    
    \item \texttt{lis3dh\_read\_int\_source}: Reads the interrupt source register, 
    which also clears the interrupt.
\end{itemize}

\subsubsection{Data Reading}

\begin{itemize}
    \item \texttt{lis3dh\_read\_acc}: Reads acceleration registers, converts raw 
    values to milligravity (mg) units, and stores them in the data structure.
    
    \item \texttt{lis3dh\_read\_temp}: Reads the temperature register and converts 
    the value to degrees Celsius.
\end{itemize}

\subsubsection{Low-Level Functions}

Four low-level functions have been designed in total. The functions \\
\texttt{lis3dh\_read\_reg} and \texttt{lis3dh\_write\_reg} read and write a single 
accelerometer register, while \texttt{lis3dh\_read\_regs\_array} and 
\texttt{lis3dh\_write\_regs\_array} read and write multiple consecutive registers. 
These four functions are essentially wrappers for the HAL functions 
\texttt{HAL\_I2C\_Mem\_Write} and \texttt{HAL\_I2C\_Mem\_Read}. These four functions are the ones that the functions explained above use to access the LIS3DH registers via the I2C protocol
